%% LyX 2.5.1 created this file.  For more info, see https://www.lyx.org/.
%% Do not edit unless you really know what you are doing.
\documentclass[12pt,english]{article}
\usepackage{mathpazo}
\usepackage[T1]{fontenc}
\usepackage[latin9]{inputenc}
\usepackage{float}
\usepackage{amsmath}
\usepackage{amssymb}
\usepackage{geometry}
\geometry{verbose,tmargin=2cm,bmargin=2cm,lmargin=3cm,rmargin=3cm}
\pagestyle{empty}
\usepackage{setspace}
\usepackage[authoryear]{natbib}
\onehalfspacing

\makeatletter
\@ifundefined{date}{}{\date{}}
%%%%%%%%%%%%%%%%%%%%%%%%%%%%%% User specified LaTeX commands.
% general
\usepackage[titletoc]{appendix}
\usepackage{graphicx}
\usepackage{placeins}
\usepackage{tikz}

% algorithm
\usepackage[ruled,vlined,linesnumbered]{algorithm2e}
\IncMargin{2.5cm}
\DecMargin{2cm}
\usepackage{fullwidth}
\usepackage{enumitem}
\setlist{leftmargin=1.7cm}

% tables
\usepackage{tabularx, siunitx, multirow, booktabs}
\begingroup
% Allow `_` and `:` in macro names (LaTeX3 style)
\catcode`\_=11
\catcode`\:=11
% Internal code of `S`
\gdef\tabularxcolumn#1{%
    >{\__siunitx_table_collect_begin:Nn S{} }%
    p{#1}%  <- this is different (is `c` in normal `S`)
    <{\__siunitx_table_print:}%
}
\endgroup

% figures
\usepackage{subfig}
\usepackage{caption}
\captionsetup[subfloat]{position=top}

% footnotes
\setlength{\skip\footins}{1cm}
\usepackage[hang,splitrule]{footmisc}
\setlength{\footnotemargin}{0.3cm} %.5
\setlength{\footnotesep}{0.4cm}

% code
\usepackage{xcolor}
\usepackage{listings}

\definecolor{codegray}{rgb}{0.5,0.5,0.5}
\definecolor{background}{HTML}{F5F5F5}
\definecolor{keyword}{HTML}{4B69C6}
\definecolor{string}{HTML}{448C27}
\definecolor{comment}{HTML}{448C27}

\usepackage{inconsolata}
\lstdefinestyle{mystyle}{
    commentstyle=\color{comment},
    keywordstyle=\color{keyword},
    stringstyle=\color{string},
    basicstyle=\ttfamily,
    breakatwhitespace=false,         
    breaklines=true,                 
    captionpos=b,                    
    keepspaces=true,                                    
    numbersep=5pt,                  
    showspaces=false,                
    showstringspaces=false,
    showtabs=false,
    tabsize=4,
	showlines=true
}

\lstset{style=mystyle}

% manual
\usepackage{enumitem}
\setlist[enumerate]{leftmargin=1cm}
\setlist[itemize]{leftmargin=0.5cm}
% Added by lyx2lyx
\setlength{\parskip}{\smallskipamount}
\setlength{\parindent}{0pt}

\makeatother

\usepackage{babel}
\begin{document}
\title{{\large\textsc{- GEModelTools -}}\textsc{}\\
\textsc{Heterogeneous Agent }\\
\textsc{NeoClassical Model (HANC)}}
\author{Rapha\"el Huleux}
\maketitle

\section*{Model}

We consider a \emph{closed} economy with heterogeneous agents and
\emph{flexible prices} \emph{and wages}.

Time is discrete and indexed by $t$. There is a continuum of households
indexed $i$.

\paragraph*{Households.}

Households are heterogeneous \emph{ex ante} with respect to their
discount factor, $\beta_{i}$, and \emph{ex post }with respect to
their productivity, $z_{it}$, and assets, $a_{it-1}$. Each period
household exogenously supply $\ell_{it}=z_{it}$ units of labor, and
choose consumption $c_{it}$ subject to a no-borrowing constraint,
$a_{it}\geq0$. Households have \emph{perfect foresight} wrt. to the
interest rate and the wage rate, $\{r_{t},w_{t}\}^{\infty}_{t=0}$,
and solve the problem
\begin{align}
v_{t}(\beta_{i},z_{it},a_{it-1}) & =\max_{a_{it},c_{it}}\frac{c^{1-\sigma}_{it}}{1-\sigma}+\beta\mathbb{E}_{t}\left[v_{t+1}(\beta_{i},z_{it+1},a_{it})\right]\label{eq:Bellman}\\
 & \text{s.t.}\nonumber \\
\ell_{it} & =z_{it}\nonumber \\
a_{it}+c_{it} & =(1+r_{t})a_{it-1}+w_{it}z_{it}\nonumber \\
\log z_{it} & =\rho_{z}\log z_{it-1}+\psi_{it}\,\,\,,\psi_{it}\sim\mathcal{N}(\mu_{\psi},\sigma_{\psi}),\,\mathbb{E}[z_{it}]=1\,\nonumber \\
a_{it} & \geq0,\nonumber 
\end{align}
where implicitly $v_{t}(\beta_{i},z_{it},a_{it-1})=v(\beta_{i},z_{it},a_{it-1},\{r_{\tau},w_{\tau}\}^{\infty}_{\tau=t})$. 

We denote optimal policy functions by $a^{\ast}_{t}(\beta_{i},z_{it},a_{it-1})$,
$\ell^{\ast}_{t}(\beta_{i},z_{it},a_{it-1})$, and $c^{\ast}_{t}(\beta_{i},z_{it},a_{it-1})$.
The distribution of households \emph{before} the realization of idiosyncratic
shocks, i.e. over $\beta_{i}$, $z_{it-1}$ and $a_{it-1}$, is denoted
$\underline{\boldsymbol{D}}_{t}$. The distribution of households
\emph{after} the realization idiosyncratic shocks, i.e. over $\beta_{i}$,
$z_{it}$ and $a_{it-1}$, is denoted $\boldsymbol{D}_{t}$.

Central aggregate variables are
\begin{align}
A^{hh}_{t} & =\int a^{\ast}_{t}(\beta_{i},z_{it},a_{it-1})d\boldsymbol{D}_{t}\label{eq:A_hh}\\
 & =\boldsymbol{a}^{\ast\prime}_{t}\boldsymbol{D}_{t}\nonumber \\
L^{hh}_{t} & =\int\ell^{\ast}_{t}(\beta_{i},z_{it},a_{it-1})d\boldsymbol{D}_{t}\label{eq:L_hh}\\
 & =\boldsymbol{\ell}^{\ast\prime}_{t}\boldsymbol{D}_{t}\nonumber \\
C^{hh}_{t} & =\int c^{\ast}_{t}(\beta_{i},z_{it},a_{it-1})d\boldsymbol{D}_{t}\label{eq:C_hh}\\
 & =\boldsymbol{c}^{\ast\prime}_{t}\boldsymbol{D}_{t}.\nonumber 
\end{align}
To solve the model, we define the beginning-of-period value function
as
\begin{equation}
\underline{v}_{t}\left(\beta_{i},z_{it-1},a_{it-1}\right)=\mathbb{E}_{t}\left[v_{t}(\beta_{i},z_{it},a_{it-1})\right].\label{eq:beginning-of-period-v}
\end{equation}
The \emph{envelope condition} implies
\begin{equation}
\underline{v}_{t,a}(\beta_{i}z_{it-1},a_{it-1})=\frac{\partial\underline{v}_{t}(\beta_{i}z_{it-1},a_{it-1})}{\partial a_{it-1}}=\mathbb{E}_{t}\left[(1+r_{t})c^{-\rho}_{it}\right].\label{eq:Envelope}
\end{equation}
The \emph{first condition} for consumption implies
\begin{equation}
c^{-\rho}_{it}=\beta\underline{v}_{t+1,a}(\beta_{i}z_{it},a_{it}).\label{eq:FOC}
\end{equation}


\paragraph*{Firms. }

A representative firm rent capital, $K_{t-1}$, and hire labor, $L_{t}$,
to produce goods, with the production function
\begin{equation}
Y_{t}=\Gamma_{t}K^{\alpha}_{t-1}L^{1-\alpha}_{t},\,\,\,\alpha\in(0,1),\label{eq:Y}
\end{equation}
where $\Gamma_{t}$ is technology and considered an exogenous shock.
Capital depreciates with the rate $\delta$. Profit maximization by
\[
\max_{K_{t-1},L_{t}}Y_{t}-w_{t}L_{t}-r^{k}_{t}K_{t-1}
\]
,implies the standard pricing equations\textbf{
\begin{align}
r^{k}_{t} & =\alpha\Gamma_{t}(K_{t-1}/L_{t})^{\text{\ensuremath{\alpha}-1}}\label{eq:rk}\\
w_{t} & =(1-\alpha)\Gamma_{t}(K_{t-1}/L_{t})^{\alpha},\label{eq:w}
\end{align}
}where $r^{k}_{t}$ is rental rate of capital and $w_{t}$ is the
wage rate. 

\paragraph*{Mutual fund.}

A zero-profit mutual fund owns all the capital. It take deposits from
households, $A_{t}$, and pay a real return of $r_{t}=r^{k}_{t}-\delta$.
It balance sheet is $A_{t}=K_{t}$.

\paragraph*{Market clearing.}

Market clearing requires
\begin{align}
\text{Capital: }K_{t} & =A_{t}=A^{hh}_{t}\label{eq:clearing_A}\\
\text{Labour: }L_{t} & =L^{hh}_{t}=1\label{eq:clearing_L}\\
\text{Goods: }Y_{t} & =C^{hh}_{t}+\underset{=I_{t}}{\underbrace{K_{t}-K_{t-1}+\delta K_{t-1}},}\label{eq:clearing_Y}
\end{align}
where $I_{t}$ is investment.

\newpage{}

\section*{Equation system}
\begin{enumerate}
\item Shocks: $\boldsymbol{Z}=\{\boldsymbol{\Gamma}\}$
\item Unknowns: $\boldsymbol{U}=\{\boldsymbol{K},\boldsymbol{L}\}$
\item Targets: $\{A_{t}-A^{hh}_{t}\}$ (asset market clearing) and $\{L_{t}-L^{hh}_{t}\}$
(labor market clearing)
\item Aggregate variables: $\boldsymbol{X}=\{\boldsymbol{\boldsymbol{\Gamma}},\boldsymbol{K},\boldsymbol{r},\boldsymbol{w},\boldsymbol{L},\boldsymbol{C},\boldsymbol{Y},\boldsymbol{A},\boldsymbol{A}^{hh},\boldsymbol{C}^{hh},\boldsymbol{L}^{hh}\}$
\item Household inputs: $\boldsymbol{X}^{hh}_{t}=\{\boldsymbol{r},\boldsymbol{w}\}$
\item Household outputs: $\boldsymbol{Y}^{hh}_{t}=\{\boldsymbol{A}^{hh},\boldsymbol{C}^{hh},\boldsymbol{L}^{hh}\}$
\end{enumerate}
This implies the equation system{\small
\begin{align}
\boldsymbol{H}(\boldsymbol{K},\boldsymbol{L},\boldsymbol{\Gamma}) & =\boldsymbol{0}\Leftrightarrow\label{eq:H_Ayagari}\\
\left[\begin{array}{c}
A_{t}-A^{hh}_{t}\\
L_{t}-L^{hh}_{t}
\end{array}\right] & =\left[\begin{array}{c}
0\end{array}\right],\,\,\,\forall t\in\{0,1,\dots,T-1\},\nonumber 
\end{align}
}where we have
\begin{align*}
r_{t} & =\alpha\Gamma_{t}(K_{t-1}/L_{t})^{\alpha-1}-\delta\\
w_{t} & =(1-\alpha)\Gamma_{t}\left(\frac{r_{t}+\delta}{\alpha\Gamma_{t}}\right)^{\frac{\alpha}{\alpha-1}}\\
A_{t} & =K_{t}\\
A^{hh}_{t} & =\boldsymbol{a}^{\ast\prime}_{t}\boldsymbol{D}_{t}\\
L^{hh}_{t} & =\boldsymbol{\ell}^{\ast\prime}_{t}\boldsymbol{D}_{t}\\
\boldsymbol{D}_{t} & =\Pi^{\prime}_{z}\underline{\boldsymbol{D}}_{t}\\
\underline{\boldsymbol{D}}_{t+1} & =\Lambda_{t}\boldsymbol{D}_{t}\\
 & \boldsymbol{\underline{D}}_{0}\text{ is given}.
\end{align*}


\section*{\newpage Implementation}

The \textbf{files }are: \lstinline{block.py}, \lstinline{household_problem.py},
\lstinline{steady_state.py}, and \lstinline{HANCModel.py}.

The results are produced in the \textbf{notebook} \lstinline{HANC.ipynb}. 

The basic \textbf{model definition} is in \lstinline{HANCModel.py}:

\vspace{2mm}\begin{lstlisting}[language=Python,basicstyle=\linespread{1.3}\ttfamily\footnotesize,
numbers=left,frame=single,backgroundcolor=\color{background}]
class HANCModelClass(EconModelClass,GEModelClass):    
    def settings(self): ...
    def setup(self): ...
    def allocate(self): ...
    prepare_hh_ss = steady_state.prepare_hh_ss
    find_ss = steady_state.find_ss
\end{lstlisting}\vspace{2mm}

The \textbf{namespaces} are:
\begin{itemize}
\item \lstinline{.par}: All parameters (no $t$ subscript)
\item \lstinline{.ss} and \lstinline{.path}: All actual variables (with
$t$ subscript)\\
Note: The variables in \lstinline{blocks.py} are in \lstinline{.path}
\end{itemize}
\textbf{Step 1.} In the \lstinline{.settings} method specify household
variables, the aggregate shocks, unknowns, targets and blocks, and
a function for solving the household problem one step backwards:

\vspace{2mm}\begin{lstlisting}[language=Python,basicstyle=\linespread{1.3}\ttfamily\footnotesize,
numbers=left,frame=single,backgroundcolor=\color{background}]
def settings(self):

    # a. namespaces (typically not changed)
    self.namespaces = ['par','ini','sim','ss','path']
    
    # b. household
    self.grids_hh = ['a'] # grids
    self.pols_hh = ['a'] # policy functions
    self.inputs_hh = ['r','w'] # direct inputs
    self.inputs_hh_z = [] # transition matrix inputs
    self.outputs_hh = ['a','c','l'] # outputs
    self.intertemps_hh = ['vbeg_a'] # intertemporal variables

    # c. GE
    self.shocks = ['Gamma'] # exogenous shocks
    self.unknowns = ['K','L'] # endogenous unknowns
    self.targets = ['clearing_A','clearing_L'] # targets = 0
    self.blocks = [ # list of strings to block-functions
        'blocks.production_firm',
        'blocks.mutual_fund',
        'hh', # household block
        'blocks.market_clearing']

    # d. functions
    self.solve_hh_backwards = None
\end{lstlisting}\vspace{2mm}

\textbf{Step 2.} In the \lstinline{.setup} method set all \emph{independent}
parameters. At the minimum:

\vspace{2mm}\begin{lstlisting}[language=Python,label=lst:setup,basicstyle=\linespread{1.3}\ttfamily\footnotesize,
numbers=left,frame=single,backgroundcolor=\color{background}]
def setup(self):
    par = self.par
    par.Nfix = 3 # number of fixed types
    par.Nz = 7 # number of discrete stochastic states
    par.Na = 300 # number of asset grid points
\end{lstlisting}\vspace{2mm}

\textbf{Step 3.} In the \lstinline{.allocate} method set all \emph{dependent}
parameters. At the minimum:

\vspace{2mm}\begin{lstlisting}[language=Python,label=lst:setup,basicstyle=\linespread{1.3}\ttfamily\footnotesize,
numbers=left,frame=single,backgroundcolor=\color{background}]
def allocate(self):         
	self.allocate_GE() # allocate on aggregate variables
\end{lstlisting}\vspace{2mm}

\textbf{Step 4. }Write the \lstinline{blocks.py} file with functions
in the following format:

\vspace{2mm}\begin{lstlisting}[language=Python,label=lst:setup,basicstyle=\linespread{1.3}\ttfamily\footnotesize,
numbers=left,frame=single,backgroundcolor=\color{background}]
import numba as nb
@nb.njit
def block_name(par,ini,ss,input1,input2,...,output1,output2,...):
    output1[:] = ...
    output2[:] = ...
\end{lstlisting}
\begin{itemize}
\item \vspace{-4mm}\textbf{Note I:} The order of the function arguments
does not matter, but good practice is inputs first, and outputs last.
\item \vspace{-2mm}\textbf{Note II:} All aggregate variables in $\boldsymbol{X}$
must be set in blocks, except for the outputs of the household block,
$\boldsymbol{Y}^{hh}_{t}$.
\item \vspace{-2mm}\textbf{Check:} At this stage it is possible to check,
that the inputs and outputs of all blocks are derived correctly by
the code. The DAG can also be produced. Run:

\vspace{2mm}\begin{lstlisting}[language=Python,label=lst:setup,basicstyle=\linespread{1.3}\ttfamily\footnotesize,
numbers=left,frame=single,backgroundcolor=\color{background}]
model = HANCModelClass()
model.info()
model.draw_DAG(figsize=(10,10))
\end{lstlisting}
\item \vspace{-2mm}\textbf{Explanation: }This works because \lstinline{GEModelTools}
automatically reads the arguments of the functions \lstinline{blocks.py}
in the order determined in the \lstinline{.blocks} attribute of the
model.
\end{itemize}
\newpage\textbf{Step 5. }In the \lstinline{household_problem.py}
file write the function \lstinline{solve_hh_backwards} to solve the
household problem in the following format:\vspace{2mm}\begin{lstlisting}[language=Python,label=lst:setup,basicstyle=\linespread{1.3}\ttfamily\footnotesize,
numbers=left,frame=single,backgroundcolor=\color{background}]
@nb.njit(parallel=True)        
def solve_hh_backwards(par,z_trans,arguments)
	
	# arguments:
    #  inputs_hh+inputs_hh_z -> r,w [scalars]
    #  intertemps_hh -> vbeg_a, vbeg_a_plus [shape=(Nfix,Nfix,Na)]
	#  outputs_hh -> a,c,l [shape=(Nfix,Nfix,Na)]
	
	# content of code:
	#  given r,w and vbeg_a_plus
	#  derive outputs and vbeg_a

\end{lstlisting}

\textbf{Step 6. }In the \lstinline{steady_state.py} file write the
function \lstinline{prepare_hh_steady} to set grids, transition matrix
of stochastic discrete states, initial distribution and initial guesses
for intertemporal variables. At the minimum:\vspace{2mm}\begin{lstlisting}[language=Python,label=lst:setup,basicstyle=\linespread{1.3}\ttfamily\footnotesize,
numbers=left,frame=single,backgroundcolor=\color{background}]
def prepare_hh_ss(model):
    par = model.par
    ss = model.ss
    par.a_grid[:] = ... # shape=(Na)
    par.z_grid[:] = ... # shape=(Nz)
	ss.z_trans[:,:,:] = ... # shape=(Nz,Nz)
	ss.Dbeg[:,:,:] = ... # shape=(Nfix,Nz,Na), sum to 1
    ss.vbeg_a[:,:,:] = ... # shape=(Nfix,Nz,Na)
\end{lstlisting}
\begin{itemize}
\item \vspace{-4mm}\textbf{Note: }The \lstinline{.prepare_hh_ss} is called
internally by \lstinline{GEModelTools} \\
whenever\textbf{ }\lstinline{.solve_hh_ss} is called.
\item \textbf{Check:} At this stage it is possible to check, that the household
problem can be solved for steady state values you choose.

\vspace{2mm}\begin{lstlisting}[language=Python,label=lst:setup,basicstyle=\linespread{1.3}\ttfamily\footnotesize,
numbers=left,frame=single,backgroundcolor=\color{background}]
model.ss.r = 0.02 # arbitrary number
model.ss.w = 1.0 # arbitrary number
model.solve_hh_ss(do_print=True) # calls prepare_hh_ss
model.simulate_hh_ss(do_print=True)
\end{lstlisting}

And that the solution has converged properly:

\vspace{2mm}\begin{lstlisting}[language=Python,label=lst:setup,basicstyle=\linespread{1.3}\ttfamily\footnotesize,
numbers=left,frame=single,backgroundcolor=\color{background}]
model.test_hh_path()
\end{lstlisting}
\end{itemize}
\newpage\textbf{Step 7. }In the \lstinline{steady_state.py} file
write the function \lstinline{find_ss} to find the steady state.
This can be formulated in many ways. The structure could e.g. be:\vspace{2mm}\begin{lstlisting}[language=Python,label=lst:setup,basicstyle=\linespread{1.3}\ttfamily\footnotesize,
numbers=left,frame=single,backgroundcolor=\color{background}]
import numba as nb
def find_ss(model)
	# a. guess on some variables
	#  could be e.g. ss.x = par.x_ss.
	# b. derive some more variables analytically
	# c. solve household problems
	model.solve_hh_ss(do_print=do_print)     
	model.simulate_hh_ss(do_print=do_print)
	# d. derive more variables analytically
	# e. check remaning equations -> update par.x_ss and return to step a.?
\end{lstlisting}
\begin{itemize}
\item \vspace{-4mm}\textbf{Note: }The function \lstinline{find_ss} should
set ALL aggregate variables, $\boldsymbol{X}$, in \lstinline{.ss},
and it should not dependent on existing values, except through \lstinline{.par}.
To verify this you can call \lstinline{.set_ss_to_nan()} just before,
or in the beginning.
\item \vspace{-2mm}\textbf{Check:} At this stage it is possible to check,
that the household problem can be solved for steady state values you
choose.

\vspace{2mm}\begin{lstlisting}[language=Python,label=lst:setup,basicstyle=\linespread{1.3}\ttfamily\footnotesize,
numbers=left,frame=single,backgroundcolor=\color{background}]
model.test_ss() # check for NaN values in .ss
model.test_hh_ss() # check for proper convergence of household problem
model.test_path() # check for consistency of .ss with blocks.py
model.test_jacs() # test computation of Jacobians
\end{lstlisting}
\end{itemize}
\textbf{Solution: }The non-linear transition path can now be found
as

\vspace{2mm}\begin{lstlisting}[language=Python,label=lst:setup,basicstyle=\linespread{1.3}\ttfamily\footnotesize,
numbers=left,frame=single,backgroundcolor=\color{background}]
model.compute_jacs()
model.find_transition_path(...)
\end{lstlisting}
\end{document}
